% Abhakseite „Das kann ich“ – Zone nur im Gesamt (\mitzone)
%% TODO Ich-kann-Sätze: Platzhalter wie in den Aufgabendateien
\begin{abhakseite}
\ifdefined\mitzone
\abhakgruppe{Kennst du schon}
\abhak{1}{Erwartungswert und Standardabweichung als Kenngrößen deuten}
\abhak{2}{Flächeninhalt unter einem Graphen als Größe deuten}
\abhak{3}{Binomialverteilung und ihr Säulendiagramm lesen}
\abhak{4}{Gegenereignis und Prozentrechnung}
\abhak{5}{Den Rechner für Verteilungswahrscheinlichkeiten einsetzen}
\abhak{6}{Flächeninhalt unter einem Graphen als Größe deuten – Fehler finden}
\abhak{7}{Flächeninhalt unter einem Graphen als Größe deuten – selbst rechnen}
\fi
\abhakgruppe{Einheit 1 · Modell und Glockenkurve: die Dichtefunktion lesen und skizzieren (μ als Symmetrieachse und Maximumsstelle, σ als Breitenmaß}
\abhak{8}{Modell und Glockenkurve}
\abhak{9}{Modell und Glockenkurve – Fehler finden, Begründen}
\abhakgruppe{Einheit 2 · Einheit 2}
\abhak{10}{Wahrscheinlichkeiten berechnen}
\abhak{11}{Wahrscheinlichkeiten berechnen – Fehler finden, Begründen}
\abhakgruppe{Einheit 3 · Einheit 3}
\abhak{12}{Umkehraufgaben und Argumente}
\abhak{13}{Standardabweichung einer Normalverteilung aus einer vorgegebenen Wahrscheinlichkeit bei festem Erwartungswert ermitteln}
\abhak{14}{Umkehraufgaben und Argumente – Fehler finden, Begründen}
\end{abhakseite}
